\section{模型假设}

\begin{enumerate}
  \item \textbf{实时信息集假设。}预测原点之后的观测不进入估计、调参或阈值选择；统计序列按实际
  发布滞后截断。固定下载版本可能被后续修订，但不回填未来信息。
  \item \textbf{交易日映射假设。}非交易日事件映射到其后首个共同交易日。短窗口仍可能叠加其他
  新闻，故CAR只解释为事件相关异常；常规正态区间为描述性结果，正式显著性依赖状态匹配安慰剂。
  \item \textbf{递归结构识别假设。}月内全球供给响应最慢，全球真实活动可当期响应供给，实际油价
  可当期响应前两者。递归排序为“供给增长---全球真实活动---实际Brent收益”，并用替代排序探针
  检验结论依赖性。
  \item \textbf{局部动态假设。}主要宏观响应在12个月内体现；PPI、CPI、IAV和原油成本的被解释量
  为未来水平，汇率为相对冲击前基期的累计对数变化。移动块bootstrap保留月度序列依赖。
  \item \textbf{跨国可比边界。}六个对照国的官方零售燃油对数收益可以比较；中国调价指数代理与
  固定品零售价不可直接等同，因此不进入正式零售价格排名。面板系数只表示对照国相对中国的差异。
  \item \textbf{局部政策反事实假设。}关闭临时调控时，将官方公告核验的应调与实调差额加入中国
  调价指数代理路径，其余宏观环境保持不变。反事实用0至6阶动态核及结果变量AR(1)递推，属于条件
  情景而非一般均衡福利估计。
  \item \textbf{风险和规则约束假设。}FHS可重采样截止日前标准化残差；SAPR只在三状态、六个月、
  离散传导率规则族内寻优。平均缺口月负担为延期成本代理，不等于财政支出；最大缺口、期末缺口、
  恢复期末缺口和单月调价幅度分别作为硬约束。
\end{enumerate}
